\chapter{The classification}
\label{ch:trichotomy}

We now assemble the classification for arbitrary bounded offspring support. First we separate the
four geometric regimes by coarse invariants. We then apply the two-law matching theorem to prove
universality in the bushy regime and within every chain class. In the chain regime, the branching
semigroup is both the source of the common presentation and the invariant that distinguishes the
classes.

\section{The regimes and their obstructions}

\begin{proposition}[Full trees]
\label{thm:bushy}
\lean{ChainClasses.fullTree_quasiIsometric_binary}
\leanok
\uses{def:qi}
Let $N\in\bbN$ and let $\cT$ be a rooted subtree of the $N$-ary tree with $v1,v2\in\cT$ for
every $v\in\cT$. Then $\cT\simeq\bbB$.
\end{proposition}

\begin{proof}
\leanok
\uses{def:qi}
Put $L=\ceil{\log_2N}$. Encode each of the $N$ child indices by a binary word of length $L$,
choosing the codes of $1$ and $2$ to differ in their first bit. Replacing every edge of $\cT$
by its code path embeds $\cT$ in a binary tree $\cT'$, multiplies distances by $L$, and leaves
every point of $\cT'$ within distance $L$ of a coded vertex. Hence $\cT\simeq\cT'$.

Every coded vertex of $\cT'$ has two descendant branches because $v1,v2\in\cT$, and a vertex
inside a code path reaches such a split within $L$ steps. Map $\bbB$ recursively to the first
split below the root and then to the first splits below the two children of each image. This map
replaces every binary edge by a path of length between $1$ and $L+1$, and every vertex of
$\cT'$ lies within distance $L$ of its image. It is therefore an $(L+1)$-quasi-isometry, so
$\cT\simeq\cT'\simeq\bbB$.
\end{proof}

\begin{lemma}[General-support regime obstructions]
\label{thm:regime-obstructions-general}
\lean{ChainClasses.threeRays_ae, ChainClasses.nearLine_gSample_ae,
ChainClasses.unbounded_hairs_gSample_ae}
\leanok
Let $\theta$ be finitely supported and supercritical, and let $\cT$ be a Galton--Watson tree
with law $\theta$ conditioned on infinite diameter. Almost surely $\cT$ carries three rays
meeting only at their common initial vertex. In the chain regime every vertex lies within
bounded distance of a line, while in the bushy regime $\cT$ has hairs of unbounded depth.
\end{lemma}

\begin{proof}
\leanok
\uses{thm:regime-obstructions, thm:hair-separation, thm:three-rays,
thm:harris-general}
Write $\cT^*$ for $\cT$ when $\theta_0=0$ and for its reduced skeleton otherwise. Under the
conditioning, $\cT^*$ never dies out and is an isometric subtree of $\cT$. If its reduced law is
not concentrated at one, independence along every infinite descent implies that a split occurs
below every vertex almost surely. Choose a split below a split. Two rays descend from distinct
children of the lower split, while a third returns towards the upper split and descends through a
different child. This proves the three-ray assertion.

In the chain regime the tree has no leaves. Below the first split, every vertex lies on a line:
one half follows an infinite descendant ray, while the other climbs to an ancestral split and
leaves through a different infinite child. The finite initial stem is within the depth of the
first split of such a line, giving the uniform bound. In the bushy regime, the Harris decomposition
of \cref{thm:harris-general} attaches independent conjugate-law bushes along the infinite
skeleton. For every depth $r$, such a bush has positive probability of containing a path of
length $r$ and becoming extinct. Independent trials at infinitely many skeleton vertices
therefore produce hairs of depth at least $r$ almost surely. Intersecting over $r$ gives
unbounded hairs.
\end{proof}

\begin{proposition}[Separation within the chain regime]
\label{thm:chain-separation}
\lean{ChainClasses.chainSeparation}
\leanok
\uses{def:branching-semigroup}
Let $\theta,\theta'$ be finitely supported supercritical laws in the chain regime with
$\Lambda_{\theta}\neq\Lambda_{\theta'}$, and let $\cT,\cT'$ be independent Galton--Watson trees
with these laws. Then almost surely $\cT\not\simeq\cT'$.
\end{proposition}

\begin{proof}
\leanok
\uses{thm:reachable-arity}
After exchanging the laws if necessary, fix
$m\in\Lambda_\theta\setminus\Lambda_{\theta'}$ and set $n=m+2$. Every quasi-isometry is a
$D$-quasi-isometry for some integer $D\geq1$, so it is enough to rule out each $D$ on an event of
probability one.

First consider a chain-regime tree with law $\theta'$. If a finite connected set has total
branching excess $s=\sum_v(\deg^+(v)-1)$, then a sphere just outside it has cardinality $s+2$.
Every summand is a supported shift, so $s\in\Lambda_{\theta'}$. Thus no region in the second tree
can have the sphere count $n=m+2$.

In the first tree, write $m$ as a sum of supported shifts. With positive probability, the
corresponding finite stack of split vertices occurs with long unary chains along all its incoming
and outgoing directions. Independence along a fixed ray and the second Borel--Cantelli lemma show
that, for every length $L$, two such $L$-isolated stacks occur arbitrarily far apart almost surely.
Each stack has exactly $n$ boundary chains.

Suppose a $D$-quasi-isometry existed. Choose $L$ much larger than its distortion constants and
apply tree stability to the $n$ long boundary chains of one stack. Their images determine a sphere
in the second tree with exactly $n$ outward components. Repeating the argument with a second,
distant stack ensures that at least one image lies far enough from the root for the sphere count
to apply. This forces $n=s+2$ with $s\in\Lambda_{\theta'}$, contradicting the choice of $m$.
Hence the event of a $D$-quasi-isometry has probability zero. Taking the countable union over $D$
proves the proposition.
\end{proof}

\begin{lemma}[Transitivity across independent realisations]
\label{thm:qi-transitive}
\lean{ChainClasses.pairQI_trans}
\leanok
\uses{def:qi}
Let $P_0,\dots,P_m$ be laws of random trees such that for each $0\leq i<m$ independent
realisations of $P_i$ and $P_{i+1}$ are almost surely quasi-isometric. Then independent
realisations of $P_0$ and $P_m$ are almost surely quasi-isometric.
\end{lemma}

\begin{proof}
\leanok
Realise $\cT_0,\dots,\cT_m$ independently on one product space. Each consecutive pair has the
product law in the corresponding hypothesis, so
$\bbP(\cT_i\simeq\cT_{i+1})=1$. The intersection of these finitely many events also has
probability one. On that intersection, compose the consecutive quasi-isometries. A finite
composition is again a quasi-isometry, and it maps $\cT_0$ to $\cT_m$.
\end{proof}

\section{Universality in the bushy regime}

\begin{proposition}[Bushy cross-law product form]
\label{thm:hairy-cross}
\lean{ChainClasses.hairyCross_chart}
\leanok
Fix $D\geq D_5$ and label the two reduced skeletons by the transported classes against
independent uniform families. Then both labelled skeletons carry the product law
$\mu_D\otimes\widetilde\nu$ and $\mu_D\otimes\widetilde\nu'$ with the same class law $\mu_D$,
whose one-vertex class has mass at least $\tfrac12$.
\end{proposition}

\begin{proof}
\leanok
\uses{thm:cross-relabel, thm:product-form, thm:conditional-iid}
Condition on the arity field of either reduced skeleton. By \cref{thm:conditional-iid}, its
shapes are independent and the shape at $w$ has law $\mu_{k(w)}$. The cross-law maps of
\cref{thm:cross-relabel}, together with the independent uniform variables, push every conditional
law on both sides to the same law $\mu_D$. Hence the transported labels are conditionally
independent with a distribution that does not depend on the arities. Averaging over the arity
field gives independence between label and arity at each vertex and yields the two stated product
laws. The lower bound at the one-vertex class is part of the relabelling construction.
\end{proof}

\begin{proposition}[Universality in the bushy regime]
\label{thm:hairy-general}
\lean{ChainClasses.hairy_general_ae}
\leanok
\uses{thm:hairy-cross, thm:composite-matching, thm:glued-transfer}
Let $\theta,\theta'$ be finitely supported supercritical laws in the bushy regime, and let
$\cT,\cT'$ be independent Galton--Watson trees with these laws conditioned on infinite
diameter. Then almost surely $\cT\simeq\cT'$.
\end{proposition}

\begin{proof}
\leanok
\uses{thm:hairy-cross, thm:composite-matching, thm:glued-transfer, thm:merge,
thm:cross-relabel, thm:qi-transitive}
Use \cref{thm:hairy-cross} to label both reduced skeletons on one class graph. Encode an arity in
the common core by the balanced binary cascade of \cref{thm:composite-matching}. Encode an arity
outside the core by its finite marker chain. The branching property and the product forms identify
the resulting finite-height laws with the two process laws used by the matching engine, including
all marker coordinates.

When the supports are close enough for a single common-core comparison, the potential estimate in
\cref{thm:cross-relabel} and \cref{thm:composite-matching} give a matching automorphism with
failure probability at most $Ke^{-cD^2}$. For arbitrary bounded supports, insert a finite chain
of intermediate arity laws and use \cref{thm:qi-transitive}. A matched pair of class labels gives
a marked quasi-isometry between the corresponding shapes. The binary encoding changes only the
bounded cascade blocks, and \cref{thm:merge} and \cref{thm:glued-transfer} assemble the local maps
into a quasi-isometry of the original trees. Letting $D\to\infty$ gives the almost sure conclusion.
\end{proof}

\section{Universality within a chain class}

\begin{lemma}[Geometric finite-phase resolvent]
\label{thm:renewal-resolvent}
\lean{ChainClasses.geometric_decay_of_block_contraction}
\leanok
Let $\cD$ be the nonzero lag phases and $\cK$ the killed transition matrix of the lagging-side
process, so that a transition reaching zero contributes no mass. With $H$ and $q$ the constants
of \cref{thm:common-renewal}, $\cK^{nH}\mathbf1\leq(1-q)^n\mathbf1$ for $n\geq0$. Consequently
the Green vector $F=\sum_{m\geq0}\cK^m\mathbf1$ is finite, satisfies $\mathbf1+\cK F=F$, and
obeys $F\leq(H/q)\mathbf1$.
\end{lemma}

\begin{proof}
\leanok
\uses{thm:common-renewal}
For $z\in\cD$, the $z$-coordinate of $\cK^m\mathbf1$ is the probability that the lag process,
started at $z$, avoids zero during its next $m$ advances. By the correction-word construction in
\cref{thm:common-renewal}, conditional on any current nonzero phase, the process reaches zero
within the next $H$ advances with probability at least $q$. Iterating over blocks gives
$\cK^{nH}\mathbf1\leq(1-q)^n\mathbf1$.

Every $m$ belongs to a block of length $H$, so summing the block estimate gives
$F\leq(H/q)\mathbf1$. Fubini's theorem for nonnegative sums allows the index shift
\[
 \cK F=\sum_{m\geq1}\cK^m\mathbf1=F-\mathbf1,
\]
which is the resolvent identity.
\end{proof}

\begin{proposition}[The weighted-screen gap]
\label{thm:renewal-weight-gap}
\lean{GraphMarkovMatching.screened_uniform_bound_resolvent}
\leanok
Suppose the screen row has the form $E_{h+1}\leq g(\Psi_h,\lVert E_h\rVert)+WE_h$. If there are
scalar bounds $a_m$ with $W^m\mathbf1\leq a_m\mathbf1$ and $\sum_ma_m<\infty$, then the Green
vector $F_W=\sum_{m\geq0}W^m(u\mathbf1)$ satisfies $u\mathbf1+WF_W=F_W$ and the invariant
induction closes as in \cref{thm:invariant}. A geometric tail for a different killed kernel does
not imply this hypothesis: already in one dimension $\cK=(3/4)$ has a geometric Green function
while $W=2\cK=(3/2)$ admits no finite nonnegative $F$.
\end{proposition}

\begin{proof}
\leanok
\uses{thm:invariant, thm:renewal-resolvent}
The power bound gives
$F_W\leq(\sum_m a_m)u\mathbf1<\infty$. By Fubini's theorem for nonnegative sums,
\[
 WF_W=\sum_{m\geq1}W^m(u\mathbf1)=F_W-u\mathbf1.
\]
Thus $u\mathbf1+WF_W=F_W$, and the same monotone induction as in \cref{thm:invariant} keeps every
screen vector below $F_W$. This argument depends on powers of the weighted operator $W$, not on
the transition kernel $\cK$. Indeed, for $\cK=3/4$ the killed process has a geometric Green
function, whereas $W=2\cK=3/2$ has divergent powers and admits no finite nonnegative invariant
vector. The renewal tail alone therefore does not close the screen recursion.
\end{proof}

\begin{proposition}[Chain cross-law product form]
\label{thm:chain-cross}
\lean{ChainClasses.Matched.chainCross_qi_scale}
\leanok
Let $\theta,\theta'$ be chain-regime laws with $\Lambda_{\theta}=\Lambda_{\theta'}=\Lambda$.
Then both presented skeletons carry the product law $\mu_*\otimes\hat\nu$ and
$\mu_*\otimes\hat\nu'$ with the same presented class law $\mu_*$, the quantised law at ratio
$\theta_1$ through the coupled classes of \cref{thm:chain-coupling}, whose zero class has mass
at least $\tfrac12$. Matched presented blobs are comparable at an explicit scale.
\end{proposition}

\begin{proof}
\leanok
\uses{thm:matched-presentation, thm:chain-coupling, thm:cluster-flat}
By \cref{thm:matched-presentation}, the blob rule uses only absorbed data. The branching property
therefore makes the presented pairs i.i.d. The neck coordinate remains geometric, with ratios
$\theta_1$ and $\theta_1'$ on the two sides, and is independent of the presented arity. Apply the
randomised cuts of \cref{thm:chain-coupling} to the two neck coordinates. Their class labels now
have one common law $\mu_*$, whose zero class has mass at least $1/2$. The matched presentation
gives the arity coordinates their common law. Independence of the two coordinates yields the two
product forms. Finally, \cref{thm:cluster-flat} and the bounded renewal window compare blobs with
matched labels at the stated scale.
\end{proof}

\begin{proposition}[Universality within a chain class]
\label{thm:chain-general}
\lean{ChainClasses.exists_engine_match_chain}
\leanok
\uses{thm:chain-cross, thm:composite-matching, thm:renewal-weight-gap}
Let $\theta,\theta'$ be chain-regime laws with $\Lambda_{\theta}=\Lambda_{\theta'}$, and let
$\cT,\cT'$ be independent Galton--Watson trees with these laws conditioned on infinite
diameter. Then almost surely $\cT\simeq\cT'$.
\end{proposition}

\begin{proof}
\leanok
\uses{thm:chain-cross, thm:composite-matching, thm:matched-presentation, thm:glued-transfer}
The matched presentation of \cref{thm:matched-presentation} puts both laws on a common presented
arity law, while \cref{thm:chain-cross} gives their neck labels the common class law $\mu_*$. We
encode the presented arities by the composite cascades. At every height, the encoded fields have
the process laws required by \cref{thm:composite-matching}. The exceptional top window contributes
to a finite tilt budget, and this budget is uniform in $D$ once
$\theta_1^{D^2}\leq1/2$. The engine therefore matches the two presented label fields with failure
probability tending to zero.

Each blob is a finite multi-port tree piece, and the original tree is isometric to the assembly of
these pieces over the presented skeleton. By \cref{thm:chain-cross}, matched blobs admit uniform
marked quasi-isometries. The multi-port form of \cref{thm:glued-transfer} assembles them into a
quasi-isometry of the two original trees. Taking the union over $D$ gives probability one.
\end{proof}

\section{Proof of the classification}

\begin{proof}[Proof of \cref{thm:trichotomy}]
\leanok
\proves{thm:trichotomy}
\uses{thm:bushy, thm:regime-obstructions-general, thm:chain-separation, thm:qi-transitive,
thm:hairy-general, thm:chain-general, thm:converse, thm:bushy-semigroup, thm:reachable-arity}
We first classify the laws. If $\theta_1=1$, the realisation is the ray. Otherwise, an infinite
law with $\theta_0=\theta_1=0$ belongs to (F), a law with $\theta_0=0<\theta_1<1$ belongs to
the chain regime, and a law with $\theta_0>0$ belongs to (B). These cases are mutually exclusive
and exhaustive. In the chain case, \cref{thm:reachable-arity} associates the class
$(\mathrm C_{\Lambda_\theta})$. The bushy reduced law instead has full branching semigroup by
\cref{thm:bushy-semigroup}.

We next prove the same-class assertion. The ray case is deterministic. The full-tree case is
\cref{thm:bushy}, the bushy case is \cref{thm:hairy-general}, and the chain case at a fixed
semigroup is \cref{thm:chain-general}. Whenever an argument passes through intermediate laws,
\cref{thm:qi-transitive} composes the almost sure pairwise conclusions.

Finally, \cref{thm:regime-obstructions-general} separates the ray, full-tree, chain, and bushy
regimes by the three-ray, bottleneck, line, and hair properties. Within the chain regime,
\cref{thm:chain-separation} separates distinct branching semigroups. The elementary finite class
is disjoint from all four infinite classes because a quasi-isometry preserves boundedness. This
proves both directions of \cref{thm:trichotomy}.
\end{proof}
